\documentclass[11pt]{article}

\usepackage[margin=1in]{geometry}
\usepackage{amsmath,amssymb,amsthm}
\usepackage{graphicx}
\usepackage[hidelinks]{hyperref}

\newtheorem{theorem}{Theorem}
\newtheorem{lemma}{Lemma}

\newcommand{\Buo}{\mathsf{Buo}}
\newcommand{\uo}{\mathsf{uo}}
\newcommand{\ellp}{\ell_p}

\title{A norm-convergence partial answer for Bourgain--uo Cauchy sequences in \texorpdfstring{$\ell_p$}{ell_p}}
\author{Run \texttt{fa\_banach\_001}, lane 7}
\date{2026-06-16}

\begin{document}
\maketitle

\section*{Status}

This packet records a candidate partial result for Question 5.2 of Kania and
Swaczyna's paper \emph{Bourgain--uo sequential completeness in vector lattices}
\cite{KaniaSwaczyna}.  The source question asks whether \(\Buo\)-Cauchy
sequences converge in \(\ell_p\), \(1<p<\infty\).  We prove that they always
converge in the norm of \(\ell_p\), for every \(1\le p<\infty\).  This does
not by itself prove order convergence or \(\Buo\)-convergence, so the packet is
classified as a partial result unless the intended reading of the question is
only norm convergence.

\section*{Source Question}

The open question appears on page 12 of the source PDF.

\begin{center}
\includegraphics[width=\textwidth]{figures/open_problem_crop.png}
\end{center}

The relevant bullet asks:
\[
\text{Do \(\Buo\)-Cauchy sequences converge in \(\ell_p\) for \(1<p<\infty\)?}
\]

\section*{Definitions}

We use the source paper's definition: a sequence \((x_n)\) in a vector lattice
\(E\) is \(\Buo\)-Cauchy if, for every strictly increasing sequence
\((n_k)\), the consecutive differences
\[
        x_{n_{k+1}}-x_{n_k}
\]
converge to \(0\) in order in \(E\).

In \(\ell_p\), \(1\le p<\infty\), order convergence to \(0\) implies norm
convergence to \(0\).  This is the standard order continuity of the
\(\ell_p\)-norm.

\section*{Proof Intuition}

The \(\Buo\)-Cauchy condition is stronger than merely saying that the original
successive differences are small: it tests the successive differences of every
subsequence.  In any metric space, a sequence is Cauchy exactly when no
subsequence can be arranged to have infinitely many large consecutive jumps.
In \(\ell_p\), order-null difference sequences are norm-null, so the
\(\Buo\)-Cauchy subsequence test rules out those large metric jumps.

\section*{Main Result}

\begin{theorem}
Let \(1\le p<\infty\).  Every \(\Buo\)-Cauchy sequence in \(\ell_p\) is norm
Cauchy.  Consequently, every \(\Buo\)-Cauchy sequence in \(\ell_p\) converges
in the norm topology of \(\ell_p\).
\end{theorem}

\begin{proof}
Let \((x_n)\) be a \(\Buo\)-Cauchy sequence in \(\ell_p\).  Suppose, toward a
contradiction, that \((x_n)\) is not Cauchy in norm.  Then there is
\(\varepsilon>0\) such that for every \(N\) there are \(m,n\ge N\) with
\[
        \|x_n-x_m\|_p\ge \varepsilon .
\]
By induction we may choose strictly increasing indices
\[
        r_1<r_2<r_3<\cdots
\]
such that
\[
        \|x_{r_{2j}}-x_{r_{2j-1}}\|_p\ge \varepsilon
        \qquad (j\in\mathbb N).
\]
Indeed, after \(r_{2j-2}\) has been chosen, apply non-Cauchyness with
\(N>r_{2j-2}\) and relabel the resulting pair in increasing order.

Apply the \(\Buo\)-Cauchy hypothesis to the subsequence \((x_{r_k})\).  By
definition, the sequence of consecutive differences
\[
        d_k:=x_{r_{k+1}}-x_{r_k}
\]
converges to \(0\) in order in \(\ell_p\).  Since the \(\ell_p\)-norm is order
continuous, \(d_k\to0\) in norm.  In particular,
\[
        \|d_{2j-1}\|_p
        =\|x_{r_{2j}}-x_{r_{2j-1}}\|_p
        \longrightarrow 0,
\]
contradicting the lower bound \(\|x_{r_{2j}}-x_{r_{2j-1}}\|_p\ge\varepsilon\).
Thus \((x_n)\) is norm Cauchy.  Since \(\ell_p\) is norm complete, \((x_n)\)
converges in norm to some element of \(\ell_p\).
\end{proof}

\section*{Verification Notes}

The proof uses only two standard facts:
\begin{enumerate}
\item In a metric space, non-Cauchyness produces a subsequence with infinitely
many consecutive jumps bounded below by a fixed \(\varepsilon>0\).
\item In \(\ell_p\), \(1\le p<\infty\), order convergence implies norm
convergence, because the norm is order continuous.
\end{enumerate}
The source paper already proves a different partial fact: if a \(\Buo\)-Cauchy
sequence in \(\ell_p\) has coordinatewise limit \(x\), then \(x\in\ell_p\)
\cite[Lemma 3.5]{KaniaSwaczyna}.  The theorem above is independent of that
lemma, and gives norm convergence directly.

\section*{Limitations and Novelty Check}

This packet does not prove that the norm limit is also the order limit.  Norm
convergence in \(\ell_p\) is strictly weaker than order convergence, so the
order/\(\Buo\)-convergence version of Question 5.2 remains open here.

Cheap run indexes contained no prior packet or attempt for arXiv:2512.14949.
A bounded web search for exact phrases from the question and the source title
found the source paper and adjacent Bourgain--uo/free-lattice material, but no
obvious later resolution of the \(\ell_p\) bullet.  Because the result is a
partial norm-convergence observation, novelty confidence is moderate rather
than high.

\section*{Human Review Recommendation}

Verify the intended topology in Question 5.2.  If ``converge in \(\ell_p\)''
means norm convergence, the \(\ell_p\) bullet is answered affirmatively by the
theorem above.  If it means order convergence or \(\Buo\)-convergence, this is
a useful partial result and the remaining domination/order-boundedness step
should be reviewed separately.

\begin{thebibliography}{9}

\bibitem{KaniaSwaczyna}
T.~Kania and J.~Swaczyna,
\newblock \emph{Bourgain--uo sequential completeness in vector lattices},
\newblock arXiv:2512.14949, 2025.

\end{thebibliography}

\end{document}
